\section{Глава 8. Планирование и обучение табличными методами}
\subsection{Упражнение 8.1(с. 203)}
На рис. 8.3 непланирующий метод выглядит особенно плохо, поскольку является
одношаговым; метод с многошаговым бутстрэппингом работал бы лучше. Как вы
думаете, может ли один из методов с многошаговым бутстрэппингом, рассмотренных
в главе 7, показать такие же хорошие результаты, как метод Dyna? Объясните свой
ответ.

\begin{figure}[H]
    \centering
    \includegraphics[width=0.9\textwidth]{8.1.png}
    \caption*{Рис. 8.3. Стратегии, найденные планирующими и непланирующими агентами Dyna-Q
    к середине второго эпизода. Стрелки показывают жадные действия в каждом
    состоянии; если для некоторого состояния не показана ни одна стрелка, то
    ценности всех действий в нём одинаковы. Чёрным квадратиком обозначено
    местоположение агента.}
\end{figure}

\begin{solution}
    Нет. Ни один многошаговый метод из главы 7 не сравняется с Dyna.

    Дело в том, какие пары состояние-действие обновляет каждый метод.
    Многошаговый метод обновляет только пары на реально пройденной траектории:
    чтобы изменить ценность состояния, в него надо физически прийти ногами.
    Dyna же на каждом реальном шаге делает $n$ планирующих обновлений (здесь
    $n=50$), выбирая из модели любые ранее встреченные пары --- то есть
    обновляет и те состояния, где агента сейчас нет, не заходя в них заново.
    Поэтому на рис. 8.3 у Dyna к середине второго эпизода уже выстроена цепочка
    жадных действий через половину лабиринта, а у одношагового агента изменился
    лишь последний шаг.

    Можно возразить, что самый сильный многошаговый метод --- Монте-Карло
    ($n\to\infty$) --- в конце эпизода обновит сразу всю траекторию. Но и это
    не помогает. Во-первых, МК заперт на одной буквально пройденной траектории
    и не умеет склеивать переходы в маршруты, которые целиком не проходились.
    Во-вторых, первый эпизод длинный и сильно петляет (около $1700$ шагов), а
    при $\gamma=0.95$ отдача состояния, посещённого в начале такого пути, равна
    $\gamma^{\approx 1700}\approx 0$: сигнал от цели до него практически не
    доходит, и выученная ценность отражает не кратчайший путь, а долгие
    блуждания. Планирующие обновления Dyna вида
    $Q(s,a)\leftarrow R+\gamma\max_a Q(s',a')$ за счёт оператора $\max$
    распространяют ценность назад по кратчайшему маршруту через модель.

    Итог: многошаговый бутстрэппинг обгонит одношаговый метод, но догнать Dyna
    не сможет.
\end{solution}

\subsection{Упражнение 8.2(с. 206)}
По какой причине Dyna-агент с призом за исследование, Dyna-Q+, показывает
лучшие результаты и в первой, и во второй частях экспериментов с блокированием и коротким путём?

\begin{solution}
    Dyna-Q+ хранит для каждой пары состояние-действие счётчик $\tau$ --- число
    шагов с момента, когда её последний раз пробовали в реальной среде, --- и на
    шаге планирования прибавляет к смоделированному вознаграждению бонус
    $\kappa\sqrt{\tau}$. Принцип один на обе части: где агент давно не был,
    $\tau$ велик, бонус большой, и жадная политика сама тянет проверить это
    место. Получается направленное исследование вместо случайного.

    Первая часть, до изменения среды. В начале обучения агент потрогал лишь
    малую долю пар, а у остальных $\tau$ только растёт, поэтому бонус заставляет
    систематически пробовать непройденные действия. Так Dyna-Q+ быстро строит
    полную модель и раньше выходит на оптимальный путь. Обычный Dyna-Q
    исследует только вслепую, случайными $\varepsilon$-жадными шагами, и
    покрывает пространство медленнее --- отсюда его более низкая кривая ещё до
    всякого изменения.

    Вторая часть, после изменения. В заброшенном участке $\tau$ растёт, бонус
    надувает ценность давно не пробованного действия, ведущего туда; агент идёт
    его проверить и из реального опыта узнаёт о перемене. Здесь важна асимметрия
    двух задач. При блокировании старый путь перестаёт работать, и среда сама
    вынуждает даже обычный Dyna-Q переучиться. А при коротком пути старый
    маршрут по-прежнему ведёт к цели, давления нет: модель Dyna-Q утверждает,
    что прохода не существует, и чем дольше он планирует, тем сильнее жадная
    политика отворачивается от нужного угла. Случайно набрать целую
    последовательность шагов в ту сторону почти невозможно, поэтому короткий
    путь обычный Dyna-Q не находит никогда. Бонус Dyna-Q+ именно это и
    исправляет: он нужен, чтобы замечать не только ухудшения среды, но и её
    улучшения, которые иначе остаются незамеченными. Любопытство стоит лишних
    проверочных шагов, но в этих задачах окупается.
\end{solution}

\subsection{Упражнение 8.3(с. 206)}
Если внимательно посмотреть на рис. 8.5, то можно заметить, что различие между
агентами Dyna-Q+ и Dyna-Q немного меньше в первой части эксперимента. Чем это
можно объяснить?

\begin{figure}[H]
    \centering
    \includegraphics[width=0.75\textwidth]{fig_8_5.png}
    \caption*{Рис. 8.5. Усреднённые результаты Dyna-агентов в задаче о коротком пути в
    лабиринте. На протяжении первых 3000 шагов использовалась левая окружающая
    среда, а затем правая.}
\end{figure}

\begin{solution}
    Меньший зазор в первой части объясняется тем, что там исследованию почти
    нечего находить: среда статична, а полезный путь всего один. Преимущество
    Dyna-Q+ здесь узкое --- оно целиком за счёт более быстрого первичного
    обучения. Пока маршрут ещё не известен, бонус $\kappa\sqrt{\tau}$ заставляет
    агента систематически перебирать непройденные пары, и он покрывает лабиринт
    быстрее, чем обычный Dyna-Q со своими случайными $\varepsilon$-жадными
    шагами. Поэтому Dyna-Q+ находит оптимальный путь раньше и раньше начинает
    набирать вознаграждение оптимальным темпом --- эта фора и сдвигает его
    кривую вверх.

    Но как только путь найден, исследование перестаёт окупаться и даже начинает
    стоить: счётчик $\tau$ у давно не тронутых пар растёт, и бонус рано или
    поздно перерастает ценность оптимального действия, уводя агента на
    бессмысленный крюк. Поэтому фора сохраняется, но не увеличивается, и в
    статичной первой части агенты идут почти вровень. Разница резко растёт лишь
    во второй части, когда открывается короткий путь: тогда то же направленное
    исследование находит новый маршрут, которого обычный Dyna-Q не обнаруживает
    вовсе.
\end{solution}

\subsection{Упражнение 8.4(с. 206)}
Описанный выше приз за исследование
фактически изменяет оценки ценностей состояний и действий. Так ли это
необходимо? Предположим, что приз $\kappa\sqrt{\tau}$ используется не при
обновлении, а только при выборе действий, т.\,е. что всегда выбирается такое
действие, для которого $Q(S_t,a)+\kappa\sqrt{\tau(S_t,a)}$ достигает максимума.
Проведите эксперимент в сеточном мире, который демонстрирует сильные и слабые
стороны такого подхода.

\begin{solution}
    \textbf{Код смотри в отдельном файле репозитория.}

    Эксперимент поставлен на лабиринте с коротким путём (как в примере 8.3):
    сетка $6\times 9$, на шаге $5000$ из $10000$ справа от барьера открывается
    короткий путь, длинный при этом остаётся. Сравниваются три агента --- Dyna-Q,
    Dyna-Q+ и предложенный вариант, где приз входит только в правило выбора
    $\arg\max_a\!\big[Q(s,a)+\kappa\sqrt{\tau(s,a)}\big]$, но не в обновления.
    Параметры стандартные ($n=50$, $\alpha=1$, $\gamma=0.95$, $\varepsilon=0.1$,
    $\kappa=10^{-3}$), усреднение по $30$ запускам.

    \begin{figure}[ht]
        \centering
        \includegraphics[width=0.75\textwidth]{8_4_sol_1.png}
        \caption*{Накопленное вознаграждение. Короткий путь открывается на шаге
        $5000$ (пунктир).}
    \end{figure}

    \begin{figure}[ht]
        \centering
        \includegraphics[width=0.75\textwidth]{8_4_sol_2.png}
        \caption*{Вознаграждение за шаг (скользящее окно $500$), тот же
        эксперимент.}
    \end{figure}

    Удобнее читать график вознаграждения за шаг (второй рисунок) --- на нём видны
    оба режима.

    Сильная сторона (статичная первая фаза, до шага $5000$). Здесь исследовать
    почти нечего, и селективный вариант (зелёный) даёт самый высокий темп
    вознаграждения, а полный Dyna-Q+ (красный) --- самый низкий. Причина в том,
    что у селективного варианта $Q$ остаётся честной оценкой ценности: приз не
    вшит в неё, поэтому, найдя путь, агент чисто его эксплуатирует. У Dyna-Q+
    приз распространяется через планирование и раздувает ценности давно не
    тронутых пар, из-за чего даже жадное действие периодически оказывается
    крюком, и агент чаще сходит с оптимального маршрута. Значит, в статичной
    среде менять оценки ценностей не только не нужно, но и слегка вредно.
    (Селективный вариант обгоняет и обычный Dyna-Q: его направленное локальное
    исследование быстрее выводит на путь и не платит постоянный
    $\varepsilon$-налог чисто случайных шагов.)

    Слабая сторона (после открытия прохода на шаге $5000$). Короткий путь лежит
    далеко от используемого маршрута, и находит его только Dyna-Q+: темп
    вознаграждения подскакивает с $\approx 0.05$ до $\approx 0.08$. Селективный
    вариант остаётся на уровне длинного пути ($\approx 0.06$), как и обычный
    Dyna-Q. Это прямо вытекает из устройства приза: в селективном варианте
    $\kappa\sqrt{\tau}$ влияет только на выбор в том состоянии, где агент стоит
    сейчас, и никуда не распространяется. Поэтому он не может выстроить градиент
    ценности, который провёл бы агента через весь лабиринт к заброшенному
    участку; а так как маршрут мимо этого участка не проходит, локальный стимул
    там попросту не срабатывает, и изменение остаётся незамеченным. У Dyna-Q+
    приз, вшитый в обновления, наоборот подтягивается к состояниям-предшественникам
    и целенаправленно ведёт агента к проходу.

    Вывод. Перенос приза в правило выбора --- это размен: мы получаем честные
    оценки ценности и чистую эксплуатацию (выигрыш в статичной среде и там, где
    переучиться вынуждает сама среда), но теряем способность дотягиваться до
    далёких неисследованных областей (проигрыш, когда находку надо искать там,
    куда текущая стратегия не ходит). Так что изменять оценки ради хорошего
    поведения в общем случае не обязательно --- но именно это изменение и
    обеспечивает дальнобойное направленное исследование, без которого далёкие
    улучшения среды не обнаруживаются.
\end{solution}

\subsection{Упражнение 8.5(с. 206)}
Как можно было бы изменить табличный алгоритм Dyna-Q на стр.~201, чтобы он был
применим к стохастической окружающей среде? Почему такая модификация могла бы
показывать плохие результаты в изменяющейся окружающей среде типа рассмотренной
в этом разделе? Как можно было бы изменить алгоритм, чтобы он хорошо работал и в
стохастической, и в изменяющейся среде?

\begin{solution}
    Часть 1 (стохастическая среда). Вместо детерминированной модели
    $Model(S,A)\leftarrow R,S'$ хранить для каждой пары $(s,a)$ выборочную оценку
    распределения исходов $\hat p(s',r\mid s,a)$ --- например, счётчики
    наблюдённых пар $(s',r)$. На шаге планирования вместо единственного
    сохранённого перехода сэмплировать $(s',r)$ из этой оценки (либо делать
    ожидаемое обновление по всем исходам).

    Часть 2 (почему плохо в изменяющейся среде). Выборочная оценка усредняет всю
    историю с равными весами. После изменения среды накопленные старые
    наблюдения подавляют немногие новые, и оценка сдвигается к новой динамике
    крайне медленно. Это даже хуже детерминированного Dyna-Q: там шаг
    $Model(S,A)\leftarrow R,S'$ --- перезапись, и одно новое наблюдение мгновенно
    исправляет модель, тогда как усреднение ещё долго тянет устаревшую статистику
    и заставляет агента планировать по ней.

    Часть 3 (и стохастическая, и изменяющаяся). Нужны два независимых
    исправления, закрывающие разные беды.

    \emph{Забывание.} Оценку распределения обновлять с постоянным шагом $\beta$,
    сдвигая весь вектор вероятностей к индикатору $e_{(S',R)}$ наблюдённого
    исхода:
    \begin{equation*}
        \hat p(\cdot\mid S,A)\leftarrow (1-\beta)\,\hat p(\cdot\mid S,A)
        + \beta\, e_{(S',R)}.
    \end{equation*}
    Это выпуклая комбинация двух распределений, поэтому результат снова
    неотрицателен и суммируется в единицу (перенормировка не нужна), а старые
    наблюдения гаснут геометрически --- модель отслеживает свежую динамику (как
    постоянный шаг для нестационарной задачи из главы~2).

    \emph{Исследование.} Добавить приз Dyna-Q+: на шаге планирования прибавлять к
    награде $\kappa\sqrt{\tau(s,a)}$, где $\tau$ --- число шагов с последнего
    реального опробования $(s,a)$.

    Забывание чинит переходы, которые агент \emph{уже наблюдает заново}; приз
    гонит его вернуться к тем, которые он \emph{перестал посещать}. Поодиночке ни
    одно не справляется, вместе --- метод тянет обе сложности сразу.

    \begin{algorithm}[H]
        \DontPrintSemicolon
        \SetKwInOut{Init}{Инициализация}
        \Init{$Q(s,a)$, оценки $\hat p(\cdot\mid s,a)$, счётчики
              $\mathrm{last}(s,a)\leftarrow 0$, время $t\leftarrow 0$}
        \While{\textnormal{обучение продолжается}}{
            $S\leftarrow$ текущее состояние\;
            $A\leftarrow \varepsilon\text{-жадная}(S,Q)$\;
            Выполнить $A$; наблюдать $R,S'$\;
            $Q(S,A)\leftarrow Q(S,A)+\alpha\big[R+\gamma\max_a Q(S',a)-Q(S,A)\big]$\;
            $\hat p(\cdot\mid S,A)\leftarrow (1-\beta)\,\hat p(\cdot\mid S,A)+\beta\,e_{(S',R)}$ \tcp*{забывающая модель}
            $t\leftarrow t+1$;\quad $\mathrm{last}(S,A)\leftarrow t$\;
            \For{$n$ \textnormal{раз}}{
                $\bar S\leftarrow$ случайное ранее наблюдавшееся состояние\;
                $\bar A\leftarrow$ случайное ранее предпринятое в $\bar S$ действие\;
                $(\bar S',\bar R)\sim \hat p(\cdot\mid \bar S,\bar A)$ \tcp*{сэмпл из распределения}
                $\bar R\leftarrow \bar R+\kappa\sqrt{t-\mathrm{last}(\bar S,\bar A)}$ \tcp*{приз за исследование}
                $Q(\bar S,\bar A)\leftarrow Q(\bar S,\bar A)+\alpha\big[\bar R+\gamma\max_a Q(\bar S',a)-Q(\bar S,\bar A)\big]$\;
            }
        }
    \end{algorithm}
\end{solution}

\subsection{Упражнение 8.6(с. 213)}
В приведённом выше анализе предполагалось, что все $b$ возможных следующих
состояний равновероятны. Предположим вместо этого, что распределение сильно
несимметрично, т.\,е. некоторые из $b$ состояний возникают с гораздо большей
вероятностью, чем другие. При таком предположении увеличится или уменьшится
преимущество выборочных обновлений над полными? Аргументируйте свой ответ.

\begin{solution}
    Преимущество выборочных обновлений увеличится.

    Полное обновление суммирует по всем $b$ преемникам и стоит порядка $b$
    операций независимо от их вероятностей: на маловероятные исходы оно тратится
    наравне с частыми, хотя вклад первых в ожидание ничтожен. Скос распределения
    эту цену не снижает.

    Выборочное обновление сэмплирует $s'\sim p(\cdot\mid s,a)$, то есть само
    вкладывает вычисления пропорционально вероятности исхода (а значит, и его
    вкладу в ожидание). При сильном скосе масса собрана на немногих состояниях:
    почти все сэмплы попадают в них, разброс сэмплированной цели мал, и среднее
    приближается к истинному ожиданию
    $\sum_{s'} p(s'\mid s,a)\,[\,r+\gamma\max_a Q(s',a)\,]$ за число обновлений,
    много меньшее $b$. Маловероятные исходы при этом недооцениваются --- но они и
    должны влиять слабо, ведь их вес $p$ мал.

    Итого: при скосе выборочное обновление достигает хорошей оценки ещё быстрее,
    тогда как полное по-прежнему платит полную цену $b$. Разрыв в пользу
    выборочного растёт.
\end{solution}

\subsection{Упражнение 8.7(с. 216)}
Некоторые графики на рис.~8.8 имеют всплески и неровности на ранних участках,
особенно это относится к верхнему графику для случая $b=1$ и равномерного
распределения. Как это можно объяснить? Какие аспекты представленных данных
подтверждают вашу гипотезу?

\begin{figure}[H]
    \centering
    \includegraphics[width=0.95\textwidth]{fig_8_8.png}
    \caption*{Рис. 8.8. Относительная эффективность обновлений, распределённых равномерно в
    пространстве состояний, и обновлений, в которых используются имитированные
    траектории с единой стратегией. Все эпизоды начинаются в одном и том же
    состоянии. Результаты соответствуют случайно сгенерированным задачам двух
    разных размеров с различными коэффициентами ветвления $b$.}
\end{figure}

\begin{solution}
    По оси ординат отложена ценность стартового состояния при жадной стратегии,
    пересчитываемая по мере накопления полных обновлений. Эта ценность зависит
    только от состояний на жадном пути из старта. При равномерном распределении
    обновляются все пары подряд, но на ценность старта влияют лишь немногие из
    них, и срабатывают они в произвольном порядке, пока
    значения ещё не устоялись: кривая стоит ровно, а затем подскакивает, когда
    очередное релевантное звено обновилось или сменилось жадное действие. При
    $b=1$ каждое полное обновление одношаговое (цель $r+\gamma\max_a Q(s')$ по
    единственному преемнику, без усреднения), поэтому такие скачки особенно
    резкие.

    Гипотезу подтверждают два наблюдения. Во-первых, всплески монотонно слабеют с
    ростом $b$: при $b=1$ они выражены, при $b=3$ тише, при $b=10$ кривая почти
    гладкая --- это согласуется с тем, что при большем ветвлении цель усредняется
    по $b$ преемникам (каждый с весом $\sim 1/b$), и изменение одного из них
    двигает её лишь на малую долю. Во-вторых, неровности сидят только на раннем
    участке, а дальше все кривые выходят на плато: эффект переходный и гаснет по
    мере того, как ценности сходятся к истинным и обновления перестают что-либо
    менять.
\end{solution}

\subsection{Упражнение 8.8(с. 224)}
Повторите эксперимент, результаты которого
представлены в нижней части рис.~8.8, а затем попробуйте выполнить его же со
значением $b=3$. Обсудите полученные результаты.

\begin{solution}
    \textbf{Код смотри в отдельном файле репозитория.}

    Воспроизведён эксперимент для задачи с 10000 состояний: случайная MDP с
    ветвлением $b$, обрывом эпизода с вероятностью 0.1 и наградами из $N(0,1)$,
    полные обновления (8.1) с $\gamma=0.9$. Сравниваются равномерное распределение
    обновлений (перебор всех пар в перемешанном порядке) и распределение с единой
    $\varepsilon$-жадной стратегией из старта; по оси ординат --- ценность старта
    при жадной стратегии (оценка Монте-Карло), усреднение по 20 задачам. Случай
    $b=1$ повторяет нижнюю панель рис.~8.8, случай $b=3$ достроен.

    \begin{figure}[h]
        \centering
        \includegraphics[width=0.8\textwidth]{fig_8_8_sol.png}
        \caption*{Задача с 10000 состояний: равномерное распределение обновлений
        против распределения с единой стратегией, для $b=1$ и $b=3$.}
    \end{figure}

    При $b=1$ распределение с единой стратегией даёт большое и продолжительное
    преимущество: оно сразу концентрирует обновления на немногих состояниях,
    достижимых из старта, поэтому ценность старта взлетает за первые сотни
    обновлений. Равномерному же на 10000 состояний 20000 обновлений хватает лишь
    на один полный проход по всем парам, поэтому значения распространяются
    медленно, и кривая догоняет единую стратегию только к концу.

    При $b=3$ это преимущество резко сжимается: кривая с единой стратегией рано
    выходит на низкое плато, а равномерное к концу её догоняет и даже слегка
    обгоняет. Причина --- рост ветвления: при $b=3$ траектории из старта быстро
    разветвляются и охватывают гораздо более широкое множество состояний, так что
    фокусировка единой стратегии на окрестности старта размывается; вдобавок
    достижимая ценность ниже (плато $b=3$ заметно ниже, чем у $b=1$).

    Вывод согласуется с общим утверждением раздела: выигрыш распределения с единой
    стратегией максимален, когда состояний много, а ветвление мало, и с ростом
    $b$ он быстро тает.
\end{solution}
